\section{Artificial Intelligence Without Artificial Persons}
\label{sec:personhood}

The phrase ``artificial person'' sounds like a prediction about technical capacity. In law, it is a statement about institutional design. A legal person exists because a legal system makes an entity the point at which claims, duties, powers, liabilities, procedures, and assets converge. Fluency and autonomy may prompt a legislature to consider creating such a point. They do not create it by themselves.

This Part makes that proposition precise without denying the seriousness of the competing view. Arguments for artificial personhood correctly perceive that advanced systems can initiate consequential sequences that no individual selected step by step. They also identify a genuine demand for stable attribution. Their mistake is to treat operational performance as though it could supply the normative and remedial architecture that personhood requires. It cannot. The more direct solution is to attribute deployment to existing persons and institutions while preserving the legislature's power to create a different status expressly.

\subsection{Legal Personality Is Not a Behavioral Test}

Legal personality begins with positions, not traits. Hohfeld's analysis remains useful because it disaggregates the word ``right'' into relations with distinct correlatives and opposites.\lrfootnote{See Wesley Newcomb Hohfeld, \artcite[30--58]{hohfeld1913}. Hohfeldian ``liability'' means exposure to another's legal power, not necessarily tort liability or capacity for punishment. The table is diagnostic; it does not imply that every legal person personally exercises every incident. Its remedy row is this Article's enforcement-oriented extension of the four Hohfeldian pairs.} Adapted to an AI dispute, the structure exposes what generated performance leaves unresolved:

\begin{table}[htbp]
\centering
\caption{Normative Positions and Artificial Performance}
\label{tab:hohfeld}
\small
\begin{tabularx}{\textwidth}{>{\raggedright\arraybackslash}p{0.15\textwidth} >{\raggedright\arraybackslash}p{0.16\textwidth} >{\raggedright\arraybackslash}p{0.25\textwidth} X}
\toprule
Position & Correlative & Legal operation & What performance alone cannot establish \\
\midrule
Claim & Duty & One subject can demand conduct from another. & That the system owns the claim or can release, assign, and enforce it. \\
Privilege & No-right & One subject may act free of another's claim to the contrary. & That the system is the beneficiary of the legal permission. \\
Power & Liability & One subject can alter a legal relation by a recognized act. & That generated assent, notice, waiver, or transfer is the system's juridical act. \\
Immunity & Disability & One subject is protected from another's attempted alteration. & That the system holds a constitutional, statutory, or private-law immunity. \\
Remedy & Exposure & Institutions can order relief against assets, conduct, or status. & What property answers, who appears, and how judgment can be enforced or sanction experienced. \\
\bottomrule
\end{tabularx}
\end{table}

A model can produce factual effects at every row. It can classify an applicant, send a notice, trigger a payment, recommend a denial, or utter words that resemble a waiver. But a factual effect and a normative position are different objects of inquiry. The payment may alter the parties' accounts because a bank, customer, or statute authorized an electronic process. The notice may be effective because the law attributes a configured communication to an organization. The denial may be unlawful because a public agency or employer used the system. In each instance, legal rules supply the subject, the attribution, and the consequence.

This is why competence cannot be the criterion of legal personality. Law recognizes newborn children and profoundly cognitively disabled people as persons without making linguistic, economic, or planning performance a threshold for human status. It also creates juridical persons that have no consciousness at all. These examples establish that sentience and legal personality are neither necessary nor sufficient for one another. They do not establish that every entity displaying a person's outward behavior becomes a person.

The corporation is the critical comparison. A corporation can own property, sue and be sued, enter contracts, commit torts and crimes, survive changes in membership, and act through officers and agents. Those capacities do not arise because the corporation passed an intelligence test. Statutes and common law constitute the entity; identify its organs; attribute acts and knowledge; partition assets; specify governance, capital, and insolvency; authorize service of process; and make remedies effective.\lrfootnote{See \casecite[636]{dartmouth1819}; \casecite[492--95]{newYorkCentral1909}; \artcite{deweyCorporatePersonality1926}; \bookcite{kurkiLegalPersonhood2019}.} Corporate personhood is thus evidence for the institutional character of personality, not a behavioral precedent for machines.

A legislature could build a comparable architecture for an AI-related entity. It could require registration, dedicated capital or insurance, a representative for service, auditable governance, ownership rules, a procedure for incapacity or termination, and a specification of whose acts and knowledge count. It might create a narrow fund-like entity for a defined activity rather than a general person. But the legal work would be done by that architecture. The system's ability to converse or plan would at most provide a policy reason for enactment.

The distinction prevents a circular argument. If a chatbot is said to possess a legal power because it can produce promise-like language, its capacity is inferred from the act whose legal effectiveness is in dispute. If it is said to bear liability because its output caused harm, the conclusion assumes that causal contribution selects the legal obligor. Ordinary law does neither. A child can cause harm without full tort capacity; a machine can cause harm without any capacity; a corporation can be liable through attribution rules even though no single employee possesses every element. Status tells law how to organize those facts. Facts do not silently enact status.

\subsection{The Strongest Case for Artificial Personhood}

The personhood literature is more sophisticated than the claim that humanlike conversation makes a machine human. Lawrence Solum's foundational inquiry treated AI personhood as a genuine legal question and tested objections grounded in intelligence, judgment, and access to human experience.\lrfootnote{See \artcite{solumLegalPersonhood1992}.} Samir Chopra and Laurence White later developed a functional account in which legal systems might recognize artificial agents where doing so serves the doctrines in which they act.\lrfootnote{See \bookcite{chopraWhite2011}.} David Gunkel and other theorists press the moral and political significance of excluding entities that may become participants in social life.\lrfootnote{See \bookcite{gunkelRobotRights2018}.} Punishment and limited-status proposals likewise ask whether recognizing a machine defendant could close responsibility gaps or serve specified institutional ends.\lrfootnote{See \artcite{abbottSarchPunishingAI2019}; \artcite{chestermanAIChallenge2020}.}

Business-entity scholarship supplies a particularly hard case for the invariant. Existing organizational law may allow an algorithm to control a limited liability company and thereby exercise many private-law functions through an already constituted juridical entity.\lrfootnote{See \artcite{bayernAutonomousSystems2015}; \artcite[887--903]{lopuckiAlgorithmicEntities2018}.} That possibility proves the institutional point rather than defeating it. The company holds the claims, assets, and procedural capacity because organizational law created the company; software control of its decisions does not silently make the software a second legal person. The arrangement also exposes why governance, disclosure, capitalization, and control rules matter more than behavioral resemblance.

An especially important responsibility proposal uses \term{operational agency} while expressly declining to confer legal personhood. On this view, an entity can be analyzed as an agent for governance because it pursues objectives, responds to constraints, and selects actions in an environment even if its metaphysical or moral agency remains disputed.\lrfootnote{See \artcite{mukherjeeChangOperationalAgency2026}.} The concept accurately describes a practical feature of current systems. A tool-using model can decompose a task, observe a result, revise a plan, and act again. A legal framework that ignores that sequence will miss the source of risk.

Operational agency is therefore a useful layer but not a complete adjudicative status rule. It identifies a pattern of causally organized action and helps trace culpability. Legal agency in the private-law sense additionally requires a relation among principal, agent, and third party, including assent, acting on another's behalf, and a right of control.\lrfootnote{See \bluecite{restatementAgency2006}.} Criminal responsibility requires legally specified fault and a sanction that can operate upon the defendant. Contractual capacity requires rules about authority, consent, avoidance, and enforcement. A single operational description cannot supply those normative relations---a limit the operational-agency account itself recognizes.

Indeed, personhood can worsen the responsibility problem it is meant to solve. A nominal machine defendant with no independent assets, no welfare affected by sanction, and no governance capable of responding would function as a liability sink. Capitalization and insurance could make the node useful, but the law can require both from providers or deploying enterprises without treating the model as the primary wrongdoer. Registration can improve traceability without personhood. Strict or enterprise liability can address evidentiary gaps without imagining machine fault. Electronic transaction rules can attribute acts without turning the program into a contracting party.

The European Parliament's nonbinding 2017 resolution recommended that the Commission consider a specific ``electronic person'' status for the most sophisticated autonomous robots. The European Economic and Social Committee officially opposed such personality because it could weaken liability's preventive and remedial effects and create moral hazard; a separate open letter pressed a similar accountability objection and argued that existing entities or insurance could perform useful allocation.\lrfootnote{See \bluecite{euRoboticsResolution2017} (para. 59(f)); \bluecite{eescAI2017} (paras. 1.12, 3.33); \bluecite{openLetterElectronicPerson2018}.} The debate did not prove that electronic personality is conceptually impossible. It showed that the operative question is institutional: which claims, assets, governance rules, and liabilities would the status create, and who benefits from placing them there?

Limited personhood therefore remains a legislative design option, not a judicial inference from capability. A defensible enactment would need to state at least: the class of systems covered; the incidents of status; the representative or organ authorized to act; the property or insurance available for remedies; the attribution of knowledge and conduct; the relation to provider, deployer, owner, and user liability; disclosure to affected persons; and the procedures for modification and dissolution. Without those components, the phrase merely renames the gap.

\subsection{Bracketing the Machine Mind}

Courts need not decide whether advanced AI is or may become conscious in order to decide present disputes. That is not because consciousness is unimportant. If credible evidence of sentience or morally relevant experience emerged, it could transform ethical obligations and create powerful reasons for protective legislation. The scientific and philosophical inquiry is active, and serious scholars disagree about both possibility and evidence.\lrfootnote{See \artcite{chalmersLLMConsciousness2023}.} The judicial point is one of decisional sequence: a court can identify the current legal subjects before resolving the system's inner life.

Moral patienthood and legal personhood must therefore be separated. Moral patienthood asks whether an entity's interests or experiences deserve moral consideration. Legal personhood asks whether a legal system constitutes the entity as the holder of specified normative positions. Law can protect interests associated with animals, cultural objects, future generations, and unincorporated environmental goods by placing duties and claims in existing human or organizational subjects, without generally constituting the protected object as a person. It can also grant juridical personality for instrumental reasons unrelated to consciousness. A conclusion about one axis cannot be imported onto the other.

Rule-following presents a related temptation. A model can produce an answer that conforms to a rule, explain the rule, and criticize another answer. Yet the philosophy of normativity has long distinguished regular behavior from occupying the social status of being committed or entitled to a norm. Wittgenstein's rule-following discussion rejects the idea that a finite sign or internal interpretation mechanically fixes every future application.\lrfootnote{See \bookcite[\S\S 138--242]{wittgensteinInvestigations1953}.} Brandom's account emphasizes practices in which participants undertake commitments, demand reasons, and attribute entitlements.\lrfootnote{See \bookcite[3--66]{brandomMakingExplicit1994}.} One need not adopt either theory wholesale to see the legal distinction: a sequence evaluated as compliant is not, without more, the bearer of the duty with which it complies.

Reinforcement learning from human feedback makes the point concrete. The training process changes a model in response to preference signals so that outputs are more likely to receive favorable evaluation.\lrfootnote{See \artcite{ouyang2022training}.} That is consequential engineering. But ``optimized to be judged acceptable under a training procedure'' and ``legally obligated to do what is right'' are different descriptions. The former concerns a learned policy and its evaluators. The latter concerns a duty holder, an authoritative norm, conditions of breach, and institutional consequences. A court should not infer the second from the success of the first.

Nor can a firm's internal constitution, model specification, or safety policy replace external judgment. Such materials can be excellent evidence of intended behavior, known risks, available safeguards, and organizational commitments. They can become contract terms or help establish a standard of care where the governing doctrine permits. But a system's self-report that it has complied with a policy is not adjudication, and the policy's use of the language of values is not a grant of public authority. Internal compliance belongs in the evidentiary record; the law determines its significance.

The bracket is deliberately symmetrical. A court need not declare that a machine has no mind for all time. It also need not treat a vivid performance as proof of one. If the legal issue is whether an existing provider made a reasonable design choice, whether a user manifested assent, or whether a public agency afforded process, the consciousness question does no necessary work. Positive law and the operative doctrine can decide the case.

\subsection{A Minimum Definition for Adjudication}

The technical definition and legal status proposition should not prove one another by stipulation. This Article therefore uses a factual working definition:

\begin{quote}
An \term{AI system} is a machine-based computational system that, with varying levels of autonomy and possible adaptiveness, infers from inputs how to generate outputs---including predictions, content, recommendations, decisions, or actions---that may influence physical or virtual environments.
\end{quote}

The definition is compatible with, but not identical to, the regulatory definitions in the EU AI Act and NIST framework.\lrfootnote{See \statcite{euAIAct2024}; \bluecite{nistAIRMF2023}.} It includes embodied and nonembodied systems, probabilistic models and hybrid systems containing deterministic components. It describes operation rather than legal status.

The independent status rule is the adjudicative invariant defended above: operational autonomy does not by itself reassign a jural position. Unless governing positive law recognizes a specified AI legal platform, every claim, duty, power, liability, and remedy necessary to judgment must terminate in a recognized natural or juridical person supported by representation, assets or insurance, and procedures for service, judgment, and enforcement. The system may supply conduct, causation, evidence, or the object of regulation; it is not thereby an independent bearer of legal responsibility.

``After deployment'' remains crucial to applying the definition. A foundation model in a repository does not determine the legal character of every later application. The same model can be deployed as an offline drafting aid, a public companion, a medical-decision component, or an agent with purchasing credentials. Configuration, tools, users, data flows, and institutional authority change the risks and the actors capable of controlling them.

Probability likewise does not become a source of free will. It describes how many present models select outputs and propagate uncertainty. It does not answer who authorized use, who knew of a hazard, or who bears a duty. A deterministic expert system, a highly predictable classifier, and a stochastic language model can occupy the same subject-status position. Parameter count, model compression, benchmark score, and whether the system is locally or remotely hosted may change foreseeable risk; none is a personhood switch.

Finally, the list of inputs prevents a prompt-centric account of causation. A user instruction is only one possible trigger. A developer's system prompt may set goals or prohibitions. Retrieved documents may alter the context. Memory may carry prior exchanges forward. A tool can return data that causes the next action. A clock, sensor, incoming message, or database change can initiate a scheduled workflow. An agent loop can generate a subgoal, invoke a tool, inspect the result, and continue without a new human command.

For legal purposes, the system is therefore an \term{automatic instrumentality}: an organized mechanism through which existing persons and institutions can produce effects without contemporaneous selection of each intermediate step. The word ``automatic'' has precedent in electronic-transactions law, where an electronic agent may act without review by an individual at the time of action while the law still attributes the transaction through existing parties.\lrfootnote{See \bluecite{ueta1999} (model act \S\S~2(6), 14); \statcite{esignElectronicAgents}.} ``Instrumentality'' does not imply simplicity or passivity. A chemical plant's safety controller and a securities-trading program are instrumentalities even though their interactions are complex and time sensitive. The term instead denies a non sequitur: causal initiative is not self-created legal status.

\subsection{From Automatic Operation to Attributed Conduct}

The automatic-instrumentality account yields a rule that will carry the rest of the Article: \emph{the absence of word-by-word human control at the instant of output does not imply the absence of human or organizational control over the mechanism that produced it}. A hospital can choose which model advises clinicians, which records it sees, how uncertainty is displayed, when a recommendation is blocked, and who can override it. A companion provider can choose personas, memory, engagement objectives, age gates, crisis detection, and exit design. A government can choose which model receives access to which data and operational tools. These choices form the relevant control structure even when no employee wrote the sentence that caused injury.

The converse matters too. Ownership of model weights is not necessarily control over a later deployment. An independent developer may materially modify open weights, remove safeguards, add tools, and release an application outside the original provider's practical reach. A user may bypass a constraint in a way that is unforeseeable or in spite of reasonable precautions. A government may compel a specific configuration. The framework should not fix responsibility at the first developer merely because that actor appears earliest in the technical supply chain. It should trace the authority, resources, information, and stopping power present in the deployment.

This account resolves the Hangzhou chatbot's apparent ``zero agency'' problem. The bot could not itself intend a contractual promise. That does not mean the interaction belonged to no legally relevant actor. The platform selected and maintained an automatic mechanism, and the court properly proceeded to ask whether the provider breached a differentiated duty of care.\lrfootnote{See \casecite{hangzhouChatbot2026}.} The contract claim and the negligence analysis did not require a fictional machine mind; they required separate attribution and duty inquiries.

The next Part shows why courts must preserve that separation when the system's functional classification changes. Part IV will then turn the proposition about control into an auditable deployment map.
